\documentclass[11pt,a4paper]{article}
\usepackage[margin=1in]{geometry}
\usepackage{times}
\usepackage{hyperref}
\usepackage{enumitem}
\usepackage{titlesec}
\usepackage{booktabs}
\usepackage{array}
\usepackage{multirow}

\hypersetup{colorlinks=true, linkcolor=blue, urlcolor=blue, citecolor=blue}
\titleformat{\section}{\normalfont\Large\bfseries}{\thesection}{1em}{}
\titleformat{\subsection}{\normalfont\large\bfseries}{\thesubsection}{1em}{}

\title{\textbf{GovExam Alert: A Government Exam Notification Portal} \\[4pt]
\large Project Proposal for Software Engineering -- UCS503L}
\author{
Submitted to: Dr. Paramveer Sidhu \\
Thapar Institute of Engineering and Technology \\
Computer Science and Engineering Department \\[8pt]
Bir Angad Singh\thanks{Roll No: 1024031150} \quad
Raghav Mittal\thanks{Roll No: 1024030747} \quad
Aryan Juneja\thanks{Roll No: 1024030731}
}
\date{14 August 2026}

\begin{document}
\maketitle

\section*{1 Higher-order goal \dots secondary framing}
This project aims to improve access to public-sector opportunities by reducing the
time and effort citizens spend searching for government job and exam notifications
scattered across dozens of official websites. While the broader goal is equitable
access to public employment information, the immediate engineering objective is to
translate this into a measurable, deployable web application that aggregates,
organizes, and alerts users about relevant exam notifications.

\section*{2 Time-to-value \dots primary heuristic}
\begin{itemize}[leftmargin=1.5em]
  \item We will deliver meaningful increments quickly by prototyping the core
  browse-and-search workflow first and validating it with a small group of student
  users within a short iteration window.
  \item We will prioritize fast feedback over perfection by using weekly iterations
  and targeting CI-backed automation early, to reduce release friction.
  \item Success is defined by what can be delivered and measured in the first
  iteration (a working notification listing and search feature), then improved
  based on user feedback.
\end{itemize}

\section*{3 Problem Statement}
Millions of students and job seekers in India track government exam notifications
(e.g., SSC, UPSC, State PSCs, Railways, Banking) that are published on separate,
often poorly designed departmental websites. Common pain points include:
\begin{itemize}[leftmargin=1.5em]
  \item \textbf{Fragmentation:} notifications are spread across many official
  portals, PDFs, and newspaper ads, with no single reliable source.
  \item \textbf{Missed deadlines:} candidates frequently miss application windows
  or admit-card release dates due to lack of timely reminders.
  \item \textbf{Poor discoverability:} existing government sites are hard to
  search or filter by qualification, category, or state.
  \item \textbf{Information overload:} unofficial aggregator sites mix outdated,
  duplicate, or unverified postings with genuine ones, causing confusion.
\end{itemize}
\textbf{Why this matters now (contextual relevance):} With a growing number of
students preparing for competitive government exams, even a simple, reliable,
and well-organized notification portal can noticeably reduce missed opportunities
and search effort.

\section*{4 Proposed Solution \dots Software Proposition}
\subsection*{4.1 Overview}
Build a web-based ``GovExam Alert'' platform with the following components,
matching the high-fidelity prototype designed in Figma:
\begin{itemize}[leftmargin=1.5em]
  \item \textbf{Home dashboard:} a searchable exam feed with quick-stat cards
  at the top (Total Exams, Active, Urgent, Reminders) giving users an
  at-a-glance summary.
  \item \textbf{Exam cards:} each notification is shown as a card with an
  organization tag and an urgency badge (e.g., \emph{Ending Soon}, \emph{Urgent},
  \emph{New}), post count, last date to apply, exam date, and a live
  ``days left'' countdown.
  \item \textbf{Save and Apply actions:} each card has a \emph{Save} button
  (bookmark for later) and a direct \emph{Apply} button linking to the
  official source.
  \item \textbf{Reminders tab:} a dedicated bottom-navigation tab where users
  review notifications they have bookmarked or set alerts for.
  \item \textbf{Settings:} notification preferences (push notifications,
  deadline reminders sent a fixed number of days before the last date, and
  optional email notifications), along with theme and language preferences
  (see Settings screen details below).
  \item \textbf{Admin panel (staff-facing, not shown in the citizen prototype):}
  a small team of student-admins can add, edit, and verify notifications
  sourced from official sites before they appear on the dashboard.
\end{itemize}

\subsection*{4.2 Core workflow}
\begin{enumerate}[leftmargin=1.5em]
  \item Admin adds a notification with structured fields (title, organization,
  category, post count, last date, exam date, official link) via the admin
  panel.
  \item System validates required fields, auto-computes the urgency badge
  (\emph{New}/\emph{Ending Soon}/\emph{Urgent}) from the days-left value, and
  publishes the notification to the public dashboard.
  \item Citizen/student searches or scans the dashboard, and taps \emph{Save}
  on notifications of interest, or \emph{Apply} to go directly to the
  official source.
  \item System sends a reminder alert (push and/or email, per the user's
  Settings) a configurable number of days before the application deadline.
  \item User reviews all saved/reminded exams under the dedicated
  \emph{Reminders} tab.
\end{enumerate}

\subsection*{4.3 UI/UX prototype reference}
A high-fidelity prototype was designed in Figma covering both Android
(Material Design 3) and iOS (Human Interface Guidelines) variants, so that
the responsive web build can borrow consistent, platform-appropriate visual
patterns (card-based layout, bottom navigation, floating action elements).
Key screens are summarized in Table~\ref{tab:screens}.

\begin{table}[h]
\centering
\begin{tabular}{|p{2.8cm}|p{9.7cm}|}
\hline
\textbf{Screen} & \textbf{Key elements} \\
\hline
Home dashboard & Search bar; quick-stat cards (Total, Active, Urgent,
Reminders); exam cards with organization tag, urgency badge, post count,
last date, exam date, days-left countdown, Save and Apply buttons; bottom
navigation (Home, Reminders, Settings). \\
\hline
Settings & Guest/user profile; notification toggles (Push, Deadline
Reminders -- 3 days before, Email); theme and language preferences; privacy
policy and help/support links. \\
\hline
Design system cover & Confirms dual design-system coverage (Material Design
3 for Android, Human Interface Guidelines for iOS) used as the visual
reference for the responsive web build. \\
\hline
\end{tabular}
\label{tab:screens}
\caption{Prototype screens and their key UI elements}
\end{table}

\subsection*{4.4 Operational constraints for an educational, second-year setting}
\begin{itemize}[leftmargin=1.5em]
  \item Keep the solution usable on low-to-mid range devices via responsive web
  design, since many users access it on mobile.
  \item Avoid dependence on advanced research algorithms (e.g., no ML ranking in
  the first iteration); focus on correct workflow, search, and reminders.
  \item Design for deployability using common free-tier web hosting and managed
  databases, appropriate for a two-person, second-year team.
\end{itemize}

\section*{5 Solution Approach \dots Engineering focus}
This is an introductory software engineering course project, so we focus on core
web-application correctness, usability, and a repeatable delivery pipeline rather
than research novelty.

\subsection*{5.1 Search and filtering \dots non-research}
Search will be implemented with simple, well-understood techniques:
\begin{itemize}[leftmargin=1.5em]
  \item Structured filters on organization, category, state, and qualification
  (standard database queries, not free-text ranking).
  \item A basic keyword search over title/organization fields using SQL
  \texttt{LIKE}/full-text search, without claiming state-of-the-art relevance
  ranking.
  \item Sorting by nearest deadline as the default view, since urgency matters
  most to users.
\end{itemize}

\subsection*{5.2 Web architecture \dots web-based solution}
A typical 3-tier web architecture:
\begin{itemize}[leftmargin=1.5em]
  \item \textbf{Frontend:} responsive UI for browsing, searching, and bookmarking
  notifications (built with React/HTML-CSS-JS).
  \item \textbf{Backend API:} authentication, CRUD for notifications, bookmarks,
  and reminder scheduling logic (built with Node.js/Express or Flask).
  \item \textbf{Database:} stores users, notifications, bookmarks, and reminder
  logs (e.g., PostgreSQL/MySQL/MongoDB).
\end{itemize}

\subsection*{5.3 CI/CD and fast delivery enabler}
CI/CD will be used to:
\begin{itemize}[leftmargin=1.5em]
  \item Automatically build and run tests on every change.
  \item Produce repeatable deployment artifacts to a staging environment.
  \item Enable safe and frequent releases of incremental features, e.g., adding
  the alerts module after the core listing is stable.
\end{itemize}

\section*{6 Evaluation Criterion \dots measurable and attributable}
We define success using measurable metrics tied to the core workflow.

\subsection*{6.1 Primary evaluation metric}
\textbf{Time-to-discovery (TTD):} time taken by a test user to find a relevant
notification matching a given profile (e.g., ``state government, graduate,
clerical post'') using the portal, compared to manually searching official
sites.
\begin{itemize}[leftmargin=1.5em]
  \item \textbf{Target:} median TTD $\leq$ 60 seconds during pilot testing,
  at least 3x faster than manual search.
  \item \textbf{Attribution:} measured via timed usability tasks with pilot
  users and application logs (search-to-click timestamps).
\end{itemize}

\subsection*{6.2 Secondary evaluation metrics}
\begin{table}[h]
\centering
\begin{tabular}{|p{3.2cm}|p{7cm}|p{3.5cm}|}
\hline
\textbf{Metric} & \textbf{Description} & \textbf{Target} \\
\hline
Alert reliability & Percentage of scheduled reminders successfully delivered before the relevant deadline & $\geq$ 95\% \\
\hline
Data freshness & Average delay between an official notification being published and it appearing on the portal & $\leq$ 24 hours \\
\hline
User clarity & User-reported clarity of notification details using a simple 1--5 feedback question & Average $\geq$ 4/5 \\
\hline
Reliability & Availability of staging/production for browsing and login & $\geq$ 99\% uptime in pilot \\
\hline
\end{tabular}
\caption{Secondary evaluation metrics}
\end{table}

\subsection*{6.3 Pilot validation plan \dots high-level}
\begin{itemize}[leftmargin=1.5em]
  \item Recruit 10--15 pilot users (fellow students preparing for government
  exams) and 1--2 student-admins to simulate the verification workflow.
  \item Compare time-to-discovery and satisfaction before/after within the
  iteration window using short interviews and task timing.
\end{itemize}

\section*{7 Scalability \dots with foresight}
The proposition should scale in both theoretical and practical senses.
\begin{enumerate}[leftmargin=1.5em]
  \item \textbf{Scalable (theoretically):} the system should support growth in
  the number of notifications, categories, and concurrent users by:
  \begin{itemize}
    \item decoupling frontend and backend,
    \item enabling pagination and indexing for common search queries,
    \item using a stateless API design so the backend can be horizontally
    scaled.
  \end{itemize}
  \item \textbf{Operationally deployable (practically):} the system should run
  on common web hosting:
  \begin{itemize}
    \item managed database (e.g., a free-tier cloud database),
    \item containerized or platform-hosted deployment (e.g., Render, Railway,
    Vercel),
    \item CI/CD pipeline for staging/production.
  \end{itemize}
\end{enumerate}

\section*{8 Engine availability heuristic}
A mature set of ``engines'' (tooling/services) is available to implement this as
a web-based project, appropriate for a second-year skill level:
\begin{itemize}[leftmargin=1.5em]
  \item Web framework for REST APIs and reactive frontend pages (Express/Flask
  + React).
  \item Managed relational or document database (PostgreSQL/MongoDB Atlas free
  tier).
  \item Authentication approach (token-based, e.g., JWT, or session-based).
  \item Scheduled job/cron mechanism for sending reminder alerts.
  \item Test framework for unit/integration tests (e.g., Jest/PyTest).
  \item CI runner and staging deployment target (GitHub Actions + free-tier
  host).
\end{itemize}
We will use these mature, well-documented components to focus on workflow
design, correctness, and measurement rather than inventing new infrastructure.

\section*{9 Project scope and deliverables \dots iteration-friendly}
\begin{table}[h]
\centering
\begin{tabular}{|p{2.2cm}|p{9cm}|p{2.5cm}|}
\hline
\textbf{Iteration} & \textbf{Deliverable} & \textbf{Priority} \\
\hline
\multirow{7}{2.2cm}{9.1 Initial (first iteration)}
 & Home dashboard: search bar + quick-stat cards (Total, Active, Urgent,
 Reminders) & Must-have \\
\cline{2-3}
 & Exam cards with organization tag, urgency badge, post count, dates,
 days-left countdown & Must-have \\
\cline{2-3}
 & Save (bookmark) and Apply (official link) actions on each card &
 Must-have \\
\cline{2-3}
 & Admin panel: add/edit/publish a notification & Must-have \\
\cline{2-3}
 & Basic logging and timestamps for TTD measurement & Must-have \\
\cline{2-3}
 & CI: build + backend unit tests + linting & Must-have \\
\cline{2-3}
 & CD to staging with a single command or pipeline job & Must-have \\
\hline
\multirow{4}{2.2cm}{9.2 Subsequent} 
 & Reminders tab + Settings (push, deadline, and email notification
 toggles) & Should-have \\
\cline{2-3}
 & User accounts and category preferences & Should-have \\
\cline{2-3}
 & Keyword search with basic full-text matching & Could-have \\
\cline{2-3}
 & Admin dashboard for notification counts and freshness metrics & Could-have \\
\hline
\end{tabular}
\caption{Project deliverables by iteration}
\end{table}

\section*{10 Risks and mitigations}
\begin{table}[h]
\centering
\begin{tabular}{|p{3.7cm}|p{9.5cm}|}
\hline
\textbf{Risk} & \textbf{Mitigation} \\
\hline
Data accuracy & Incorrect or outdated notifications could mislead users; mitigate by always linking to the official source and having a student-admin verification step before publishing. \\
\hline
Data entry overhead & Manually adding notifications is time-consuming for a two-person team; mitigate by keeping the admin form minimal and scoping the pilot to a small set of exam categories. \\
\hline
Evaluation measurement ambiguity & Instrument timestamps and define clear notification-status states (Draft, Published, Expired) before development begins. \\
\hline
Operational issues in pilot & Use staging early; ensure CI/CD produces repeatable deployments before the pilot begins. \\
\hline
\end{tabular}
\caption{Risks and mitigations}
\end{table}

\section*{11 Summary}
This project proposes a deployable web platform, GovExam Alert, to reduce the
time
and effort citizens spend discovering and tracking government exam and job
notifications. It meets the course criteria by emphasizing time-to-value
through rapid iterations, implementing CI/CD to support frequent safe
releases, and using measurable evaluation criteria (especially
time-to-discovery). As a web-based application project scoped for a
third-year, three-member team, it remains focused on engineering workflow,
usability, and scalability readiness rather than research-driven novelty.

\end{document}
